\section{The short answer and the exact scope}
\label{sec:scope-and-answer}

There are two meanings of ``comparable,'' and only one is currently valid.

\begin{enumerate}[leftmargin=*]
\item Two examples are \emph{mechanism-comparable} when they use enough common
      structure that they isolate different failure mechanisms.  The signed
      SOR and FISTA star examples qualify: both use the same PageRank matrix
      family and degree-weighted adjacency work, and both ask whether a
      spectrally accelerated trajectory remains local.
\item Two algorithms are \emph{performance-comparable} only after the graph,
      seed, mathematical objective, output accuracy, initialization, stopping
      rule, and charged operations have been matched or connected by a proved
      reduction.  The two existing star examples do not qualify.
\end{enumerate}

Consequently, the statement supported by the present evidence is
\begin{quote}
Signed SOR proves that acceleration and local coordinate work are compatible
on one unregularized center-star, while the FISTA leaf-star proves that one
standard momentum mechanism can destroy locality on one regularized problem.
\end{quote}
It is \emph{not} valid to conclude that SOR is generally faster than FISTA,
or that the signed-star complexity dominates the FISTA lower bound.

\subsection{What differs in the two star theorems}

The comparison table makes every mismatch explicit.

\begin{center}
\begin{tabular}{p{0.19\textwidth}p{0.35\textwidth}p{0.35\textwidth}}
\toprule
Axis & Signed-SOR star theorem & Standard-FISTA star obstruction \\
\midrule
Optimization problem & Unregularized PageRank quadratic $f$; equivalently
$Qx=b$ & Composite RPPR objective $F_\rho=f+\alpha\rho
\|D^{1/2}x\|_1$ \\
Graph and seed & Center-seeded $K_{1,m}$ & Leaf-seeded $K_{1,m}$ \\
Output target & Degree-normalized semantic PPR error
$\varepsilon_{\rm ppr}$ & RPPR objective gap
$\varepsilon_{\rm obj}$ in the source theorem \\
Locality event & Repeated signed center/leaf coordinate relaxations & A
momentum overshoot activates the degree-$m$ center outside the optimal support
\\
Positive result & Exact
$O(1/(\sqrt\alpha\,\varepsilon_{\rm ppr}))$ charged work & Global
accelerated rate remains true; a conditional local bound requires confinement
\\
Negative result & Nonnegative-residual one-hop methods lose the square-root
scale on the same center-star & FISTA may pay $\Omega(m)$ although the RPPR
optimum has one-coordinate support \\
\bottomrule
\end{tabular}
\end{center}

Changing any one of these axes can change the answer.  In particular, the
leaf-star obstruction deliberately tunes $\rho$ to a tight KKT breakpoint;
the center-star SOR theorem deliberately chooses $m$ at the semantic output
scale.  These are different adversarial constructions.

\subsection{What this note adds}

This note supplies four layers, and labels them separately.

\begin{enumerate}[leftmargin=*]
\item \textbf{Proved here: an accuracy and objective bridge.}  RPPR has at
      most $\rho$ regularization bias in the project's semantic norm.  An
      objective gap has an explicit semantic conversion.  On a known positive
      face, RPPR is an ordinary PageRank principal system with a shifted
      right-hand side.
\item \textbf{Proved here: fixed-face and radial SOR generalizations.}  Optimal
      red--black SOR has an exact singular-mode rate on every bipartite
      principal face.  Spider prefixes are an explicit specialization, and a
      Chebyshev-shift argument gives a log-free semantic theorem for every
      complete equal-arm, hub-seeded spider.
\item \textbf{Proved here: a graph-uniform plain-SOR stop.}  A layered
      degree-skewed bipartite family preserves a far semantic error through
      $\Theta(1/\sqrt\alpha)$ complete source-first sweeps.  It rules out both
      a universal maximum-semantic envelope and product-scale work for this
      literal full-sweep plain-SOR primitive, even with the face supplied.
\item \textbf{Imported structural results: discovery beyond the star.}
      Exact response methods already discover and solve RPPR on spiders,
      trees, and bounded-block graphs at or below the desired product scale.
      These are not relabeled as SOR theorems.
\end{enumerate}

The graph-uniform algorithmic objective remains open, but the naive
complete-face full-sweep plain-SOR route is closed.  Exact nonsettled
face-shock orthogonality closes numerical transport under a supplied
geometric volume gate, but does not make support discovery or boundary
reporting free; the response theorems still do not cover arbitrary large
biconnected cores.

\subsection{Model and work contract}

Throughout, $Q,b,x^0$, and $F_\rho$ are the shared source-aligned objects.
The graph is finite, simple, undirected, and unweighted, with no isolated
vertices.  A row operation at $u$ costs $d_u$; a full sweep over $U$ costs
$\operatorname{vol}(U)=\sum_{u\in U}d_u$.  Repeated sweeps, support discovery,
boundary checks, state movement, and explicit output are charged whenever a
theorem claims them.  Put
\begin{equation}
 \|z\|_{D,\infty}:=\|D^{-1/2}z\|_\infty.
 \label{eq:semantic-x-norm}
\end{equation}
Then $\|\widehat x-x^0\|_{D,\infty}\leq\varepsilon_{\rm ppr}$ is exactly the
project's degree-normalized semantic PPR target after the map
$\widehat\pi=D^{1/2}\widehat x$.
